% Created 2018-05-16 Wed 11:57
% Intended LaTeX compiler: pdflatex
\documentclass[11pt]{article}
\usepackage[utf8]{inputenc}
\usepackage[T1]{fontenc}
\usepackage{graphicx}
\usepackage{grffile}
\usepackage{longtable}
\usepackage{wrapfig}
\usepackage{rotating}
\usepackage[normalem]{ulem}
\usepackage{amsmath}
\usepackage{textcomp}
\usepackage{amssymb}
\usepackage{capt-of}
\usepackage{hyperref}
\usepackage[margin=0.80in]{geometry}
\usepackage[url=false, backend=biber, style=ieee]{biblatex}
\usepackage{float}
\usepackage[super,negative]{nth}
\usepackage[capitalise]{cleveref}
\usepackage{pst-node,transparent,ragged2e}
\addbibresource{/home/nuk3/.spacemacs.d/org-files/bibliography.bib}
\DeclareFieldFormat[inproceedings]{citetitle}{\textit{#1}}
\DeclareFieldFormat[inproceedings]{title}{\textit{#1}}
\DeclareFieldFormat[misc]{citetitle}{#1}
\DeclareFieldFormat[misc]{title}{#1}
\renewcommand*{\bibpagespunct}{%
\ifentrytype{inproceedings}
{\addspace}
{\addcomma\space}}
\AtEveryCitekey{\ifuseauthor{}{\clearname{author}}}
\AtEveryBibitem{\ifuseauthor{}{\clearname{author}}}
\SpecialCoor
\author{Nyameko Lisa}
\date{26 April 2018}
\title{IPSP055 - Essential Industrial Design Law\\\medskip
\large Assignment 02 - 898979}
\hypersetup{
 pdfauthor={Nyameko Lisa},
 pdftitle={IPSP055 - Essential Industrial Design Law},
 pdfkeywords={},
 pdfsubject={},
 pdfcreator={Emacs 27.0.50 (Org mode 9.1.2)},
 pdflang={English}}
\begin{document}

\maketitle
\addvspace{110pt}
\centering{
\pnode(0.5\textwidth,-0.5\textheight){thisCenter}
\rput(thisCenter){%\transparent{0.25}
\includegraphics[width=2.7in]{/home/nuk3/course/llb/wipo-unisa/UNISACoatofArms.eps}}}

\justifying
\addvspace{110pt}
\section*{Declaration}
\label{sec:org39968e6}
I know that plagiarism is to use someone else’s work and pass it off as my own.\\
I know that plagiarism is wrong.\\
I confirm that this assignment is my own work.\\
I have acknowledged in the bibliography accompanying the assignment all the sources that I have used.\\
I have not directly copied without acknowledgement anything from the Internet or from any other source.\\
I have indicated every quotation and citation in a footnote or bracket linked to that quotation.\\
I have not allowed anyone else to copy my work and to pass it off as their own work.\\
I understand that if any unacknowledged copying whatsoever appears in my assignment I will receive zero per cent for the assignment.\\
I am aware of the UNISA policy on plagiarism and understand that disciplinary proceedings can be instituted against me by UNISA if I contravene this policy.\\
I indicate my understanding and acceptance of this declaration by
entering my name here under:
\begin{itemize}
\item Name: \textbf{Nyameko Lisa} (Student Number: \textbf{7874-909-3})
\end{itemize}

\subsection*{NOTE}
\label{sec:org5838cac}
Please note that footnotes will be denoted as \footnote{This is a footnote.} and will
appear at the bottom of the page.\\
References will be denoted by \cite{rsa93_designs_act} and will appear at the end of the document.
\newpage

\section{Analyse the possibility of design infringement. [20]}
\label{sec:org6b13170}

As per the definitions of [section 2(1)]\cite{rsa93_tm_act} and given that the
Republic is both a Paris Convention as well as a World Intellectual Property
Organisation (WIPO) Member State, [article 15(1)]\cite{wto17_trips}, [article
1(2)]\cite{wipo83_paris_conve_protect_ip}, \cite{wipo25_hague_agreement}, South African statutory design
legislation is in compliance with International Conventions and legal instruments.\\

As per the provisions of \cite{rsa78_patents_act}, patent rights seek to protect
the principle underlying an inventive process or product. In contrast the South
African Designs Act No. 195 of 1993 \cite{rsa93_designs_act}, provides provisions
for the protection of only the physical form of articles of manufacture intended
to be multiplied by industrial processes, [section
14(4)]\cite{rsa93_designs_act}.\\

As per the provisions of [section 20(1)]\cite{rsa93_designs_act}, the effect of
registration of the design in \textbf{Figure 2} shall grant the registered proprietor
in the Republic, for the [section 22(1)]\cite{rsa93_designs_act}, duration of
registration the exclusive right to exclude others from \textbf{making} or \textbf{using} any
article included in the class in which their design is registered and embodying
the registered design or a design not substantially different from the
registered design.

\subsection{Infringement}
\label{sec:org43e4950}

As per the provisions of [section 35(1)]\cite{rsa93_designs_act}, the registered
proprietor of the design in \textbf{Figure 2} may institute proceedings for
infringement against Deniz Karasahin, adesignaward.com, for the \textbf{unauthorised
manufacture} of their Osteoid design, which bears resemblance not substantially
different to the registered design as an article of the same class.

\subsection{Visual Judgement of Similarity}
\label{sec:org02d09b5}

As per \citetitle{corbett84_homecraft_steel_v_sm_hare}
\cite{corbett84_homecraft_steel_v_sm_hare}, the test for infringement is a purely
visual judgement according to the observation of the likely customer where Judge
Corbett  referred to Judge Eloff in preceding case law,
\begin{quote}
\textit{``to hold that while the eye is to be the eye of the Court, the
Court should view the design through the spectacles of the customer''}
\end{quote}
wherein the designs are viewed and considered both independently as well as
side by side,
\begin{quote}
\textit{``an article must have some special feature which appeals to the
eye and distinguishes it from other of its class... This must be decided by the
eye; and in applying this visual test the Court must bear in mind the reaction
of the likely customer. The two designs must be viewed and compared side by side
and also separately.''}
\end{quote}

\subsection{Novelty and Originality}
\label{sec:org748e1d1}

When analysing the possibility of infringement, the state of the art and the
degree of novelty and originality achieved by the registered design is also an
important factor, in reference to preceding case law, Judge Corbett stipulates,
\cite{corbett84_homecraft_steel_v_sm_hare} that,
\begin{quote}
\textit{``If only small differences seperate the registered design from
what has gone before, then eqaully small differences between the alleged
infringement and the regsitered design will be held to be sufficient to avoid
infringement...''}
\end{quote}
There are however significant differences between the two designs, in \textbf{shape} and / or
\textbf{configuration}, and there is no risk of a member of the public or likely customer
confusing the two designs, where in substantiating proceedings for design
infringement applicants have in the past relied on alleged infringement of their
registered designs as well as on the delict of passing off,
\citetitle{corbett84_homecraft_steel_v_sm_hare} \cite{corbett84_homecraft_steel_v_sm_hare},
\citetitle{vdwesthuizen02_slavepak_v_buddy} \cite{vdwesthuizen02_slavepak_v_buddy},
\citetitle{duplessis02_koninklijke_phillips_v_kenwood}
\cite{duplessis02_koninklijke_phillips_v_kenwood}.\\

As per the provisions of [section 35(3)]\cite{rsa93_designs_act}, the plaintiff
would likely seek and interdict [subsection (a)], assuming the prototype is still in the
developmental stages and has not as yet gone into production, and would not
likely pursue surrender of goods or damages [subsections (b)-(-d)]. Based on the
above arguments however, the plaintiff is not likely to succeed with any design
infringement actions within the Republic.

\subsection{Recommendation for design and patent protection}
\label{sec:org233d204}

The designer Deniz Karasahin, would be strongly recommended to not only seek
aesthetic and functional design protection, but also patent protection for their
Osteiod design. As per the definitions specified in [section
2]\cite{rsa78_patents_act} and [Article 27(1)]\cite{wto17_trips}, the Osteoid prototype
need satisfy the provisions for a patent to be granted under [section
25]\cite{rsa78_patents_act}, wherein subsection (1) of that section of the Act
stipulates that a patent may be granted for a \textbf{new} invention\footnote{Novel 3D-printed design.\label{org509d990}}, involving
an \textbf{inventive step}\footnote{Innovative low-intensity pulsed ultrasound system.\label{orga729b9b}} and has application to trade or
industry\footnote{Contributing to the medical healing process.}.

\subsubsection{Aesthetic design registration}
\label{sec:orgff51973}
As per the provisions of [section 1(1)]\cite{rsa93_designs_act}, which stipulates
that \textit{``...any design applied to any article, whether for the
pattern or the shape or the configuration or the ornamentation thereof, or for
any two or more of these purposes, and by whatever means is applied, having
features which appeal to and are judged solely by the eye, irrespective of the
aesthetic quality thereof...''}, as was stipulated by Judge Corbett in
\citetitle{corbett84_homecraft_steel_v_sm_hare}:
\begin{quote}
\textit{To qualify as a design, an article must have some special feature
which appeals to the eye and distinguishes it from others of its class.''}
\end{quote}
The \textbf{proprietor} of the Osteoid would be advised to pursue aesthetic design
registration for their \textbf{article}.

\subsubsection{Functional design registration}
\label{sec:org69c07e9}
Again as per the provisions of [section 1(1)]\cite{rsa93_designs_act}, the
definition of a function design is \textit{``...any design applied to
any article, either for the pattern or the shape or the configuration thereof,
or for any two or more of these purposes, and by whatever means it is applied,
having features which are necessitated by the function which the article to
which the design is applied, is to perform, and includes an integrated circuit
topography, a mask work and a series of mask works...''}.\\

The \textbf{proprietor}, Deniz Karasahin, is strongly recommended to lodge both
functional and aesthetic design registration applications in respect of their
\textbf{article}, as per the provisions of [section
14(1)(b)(ii)]\cite{rsa93_designs_act}, and the methods and principles of it's
construction\textsuperscript{\ref{org509d990}} [section 14(4)]\cite{rsa93_designs_act}, \emph{``that are not
commonplace in the art in question''} and the functional features of the
3D-printed medical cast that requires protection of purely functional design
features\textsuperscript{\ref{orga729b9b}}.

\printbibliography
\end{document}